% ---------------------------------------------------------------
% figures/static/static/tikz/circulatory-overview.tex
% Circulatory schematic. Single arrow weight per loop, color-coded
% legend swatches, and unified label font; the caption carries the scope note.
% ---------------------------------------------------------------
\begin{tikzpicture}[fig, scale=1.0,
    organ/.style={draw=wallgray, fill=wallgraylight, rounded corners=2pt,
                  font=\footnotesize, inner sep=3pt,
                  minimum width=22mm, minimum height=8mm, align=center},
    pump/.style={draw=bloodred!70!black, fill=bloodred!15, rounded corners=3pt,
                 font=\footnotesize\bfseries, inner sep=2.5pt,
                 minimum width=12mm, minimum height=8mm, align=center},
    arteryline/.style={->, line cap=round, line join=round,
                       line width=1.0pt, bloodred!80!black},
    veinline/.style ={->, line cap=round, line join=round, dashed,
                       line width=1.0pt, mathblue!70!black},
    pulmline/.style ={->, line cap=round, line join=round,
                       line width=1.0pt, accentTeal!75!black}]

  % HEART (center)
  \node[pump] (LV) at ( 1.0, 0)   {LV};
  \node[pump] (RV) at (-1.0, 0)   {RV};
  \node[pump] (LA) at ( 1.0, 1.3) {LA};
  \node[pump] (RA) at (-1.0, 1.3) {RA};
  \draw[bloodred!50!black, line width=0.6pt, dashed]
    ($(RA.north west)+(-0.18,0.45)$) rectangle ($(LV.south east)+(0.18,-0.05)$);
  \node[font=\footnotesize\bfseries, bloodred!70!black, fill=white, inner sep=1pt]
    at (0,2.05) {Heart};

  % LUNGS (top)
  \node[organ, fill=accentTealLite] (LUNG)
       at (0,3.6) {Lungs (gas exchange)};

  % SYSTEMIC LOOP nodes (left column)
  \node[organ, fill=bloodred!12, draw=bloodred!60!black] (AORTA)
       at (-4.6,0) {Aorta \& systemic arteries};
  \node[organ] (CAP)
       at (-4.6,-1.7) {Organs / capillary beds};
  \node[organ, fill=mathblueLite] (VEIN)
       at (-4.6,-3.1) {Systemic veins};

  % CORONARY (manuscript focus). Placed on the right of the heart so the
  % aortic-root coronary supply is a short arrow off LV and does not
  % cross any systemic-loop edges. Highlighted via fill + thicker border.
  \node[organ, fill=bloodredlight!70, draw=bloodred!70!black, line width=0.9pt,
        minimum width=24mm, minimum height=8mm, inner sep=2pt,
        font=\scriptsize, align=center]
       (CORO) at (3.4,-0.6)
       {coronary arteries\\[-1pt](CAS site)};

  % SYSTEMIC LOOP (oxygen-rich red, oxygen-poor blue). The LV->aorta
  % ejection arcs below the heart so it never crosses RV or the
  % coronary supply.
  \draw[arteryline] (LV.south) to[out=-90, in=-30, looseness=1.1] (AORTA.east);
  \draw[arteryline] (AORTA.south) -- (CAP.north);
  \draw[veinline]   (CAP.south)   -- (VEIN.north);
  \draw[veinline]   (VEIN.east) to[bend right=22] (RA.south west);

  % Coronary supply only (short, direct). Coronary venous return is
  % implicit and omitted to keep the schematic uncluttered.
  \draw[arteryline] (LV.east) -- (CORO.west);

  % PULMONARY LOOP (around the top). Both strokes use pulmline so the
  % legend swatch matches the rendered color.
  \draw[pulmline] (RV.north west) to[bend left=22]
    node[left, font=\scriptsize, accentTeal!65!black, pos=0.55] {pulmonary art.}
    (LUNG.south west);
  \draw[pulmline] (LUNG.south east) to[bend left=22]
    node[right, font=\scriptsize, accentTeal!65!black, pos=0.45] {pulmonary v.}
    (LA.north east);

  % LEGEND (color-coded swatches). Placed in the lower-right quadrant,
  % below the coronary box. Cardiac-period parameters live in the body
  % text / caption to keep the schematic uncluttered.
  \begin{scope}[shift={(2.55,-2.3)}]
    \draw[arteryline] (0,0.45) -- (0.85,0.45);
    \node[font=\scriptsize, anchor=west, wallgray!85!black] at (0.92,0.45)
      {systemic arterial};
    \draw[veinline] (0,-0.05) -- (0.85,-0.05);
    \node[font=\scriptsize, anchor=west, wallgray!85!black] at (0.92,-0.05)
      {systemic venous};
    \draw[pulmline] (0,-0.55) -- (0.85,-0.55);
    \node[font=\scriptsize, anchor=west, wallgray!85!black] at (0.92,-0.55)
      {pulmonary loop};
  \end{scope}
\end{tikzpicture}
